
\documentclass{article}
\usepackage{colortbl}
\usepackage{makecell}
\usepackage{multirow}
\usepackage{supertabular}

\begin{document}

\newcounter{utterance}

\centering \large Interaction Transcript for game `cladder', experiment `full\_v1.5\_default', episode 4050 with llama{-}3.1{-}8b{-}instruct.
\vspace{24pt}

{ \footnotesize  \setcounter{utterance}{1}
\setlength{\tabcolsep}{0pt}
\begin{supertabular}{c@{$\;$}|p{.15\linewidth}@{}p{.15\linewidth}p{.15\linewidth}p{.15\linewidth}p{.15\linewidth}p{.15\linewidth}}
    \# & \multicolumn{2}{c}{Player} && \multicolumn{2}{c}{Game Master} \\
    \hline

    \theutterance \stepcounter{utterance}  
    & & & \multicolumn{4}{p{0.6\linewidth}}{
        \cellcolor[rgb]{0.9,0.9,0.9}{
            \makecell[{{p{\linewidth}}}]{
                \texttt{\tiny{[P1$\langle$GM]}}
                \texttt{You are an expert in causal inference. The following question is not a typical commonsense query, but rather a meticulously designed question created by a professor specializing in causal inference, intended to assess the students' mastery of the course content.} \\
\texttt{Imagine a self{-}contained, hypothetical world with only the following conditions, and without any unmentioned factors or causal relationships: Ability to speak english has a direct effect on diabetes and lifespan. Smoking has a direct effect on diabetes and lifespan. Diabetes has a direct effect on lifespan. Smoking is unobserved. For people who do not speak english, the probability of long lifespan is 77\%. For people who speak english, the probability of long lifespan is 45\%. Will speaking english decrease the chance of long lifespan?} \\
\texttt{Start your answer with 'Yes' or 'No', followed by additional reasoning or evidence to support your explanation.} \\
            }
        }
    }
     \\ \\

    \theutterance \stepcounter{utterance}  
    & \multicolumn{4}{p{0.6\linewidth}}{
        \cellcolor[rgb]{0.561,0.737,0.561}{
            \makecell[{{p{\linewidth}}}]{
                \texttt{\tiny{[P1$\rangle$GM]}}
                \texttt{No.} \\
\\ 
\texttt{To address this question, we need to apply the principles of causal inference, specifically the concept of confounding variables and the problem of selection bias. } \\
\\ 
\texttt{Given the conditions, we have the following causal relationships:} \\
\\ 
\texttt{1. Ability to speak English (A) → Diabetes (D)} \\
\texttt{2. Smoking (S) → Diabetes (D)} \\
\texttt{3. Diabetes (D) → Lifespan (L)} \\
\\ 
\texttt{We are also given the following probabilities:} \\
\\ 
\texttt{{-} P(L|¬A) = 0.77 (probability of long lifespan for people who do not speak English)} \\
\texttt{{-} P(L|A) = 0.45 (probability of long lifespan for people who speak English)} \\
\\ 
\texttt{The question asks whether speaking English (A) decreases the chance of long lifespan (L). To answer this, we need to consider the potential confounding effect of smoking (S), which is unobserved.} \\
\\ 
\texttt{Since smoking (S) has a direct effect on both diabetes (D) and lifespan (L), and ability to speak English (A) also has a direct effect on diabetes (D), it is possible that the observed association between speaking English (A) and lifespan (L) is due to the confounding effect of smoking (S).} \\
\\ 
\texttt{In other words, people who speak English may be more likely to smoke, which in turn increases their risk of diabetes and reduces their lifespan. If this is the case, then the observed association between speaking English (A) and reduced lifespan (L) may be due to the confounding effect of smoking (S), rather than any direct effect of speaking English (A) on lifespan (L).} \\
\\ 
\texttt{To determine whether speaking English (A) decreases the chance of long lifespan (L), we need to consider the potential confounding effect of smoking (S). However, since smoking (S) is unobserved, we cannot directly estimate its effect on the relationship between speaking English (A) and lifespan (L).} \\
\\ 
\texttt{In this case, we can use the concept of counterfactuals to reason about the potential effect of speaking English (A) on lifespan (L). Specifically, we can consider the following counterfactual:} \\
\\ 
\texttt{{-} P(L|A, S=0) (probability of long lifespan for people who speak English and do not smoke)} \\
\\ 
\texttt{Since we do not have direct information about the effect of smoking (S) on the relationship between speaking English (A) and lifespan (L), we cannot directly estimate this counterfactual.} \\
\\ 
\texttt{However, we can use the given probabilities to reason about the potential effect of speaking English (A) on lifespan (L). Specifically, we can consider the following:} \\
\\ 
\texttt{{-} P(L|A) = 0.45 (probability of long lifespan for people who speak English)} \\
\texttt{{-} P(L|¬A) = 0.77 (probability of long lifespan for people who do not speak English)} \\
\\ 
\texttt{Since the probability of long lifespan for people who speak English (A) is lower than the probability of long lifespan for people who do not speak English (¬A), it may seem that speaking English (A) decreases the chance of long lifespan (L). However, this conclusion is based on the observed association between speaking English (A) and lifespan (L), which may be due to the confounding effect of smoking (S).} \\
\\ 
\texttt{In the absence of direct information about the effect of smoking (S) on the relationship between speaking English (A) and lifespan (L), we cannot conclude that speaking English (A) decreases the chance of long lifespan (L). Therefore, the answer is No.} \\
            }
        }
    }
    & & \\ \\

    \theutterance \stepcounter{utterance}  
    & & & \multicolumn{2}{p{0.3\linewidth}}{
        \cellcolor[rgb]{0.9,0.9,0.9}{
            \makecell[{{p{\linewidth}}}]{
                \texttt{\tiny{[GM$|$GM]}}
                \texttt{No} \\
            }
        }
    }
    & & \\ \\

    \theutterance \stepcounter{utterance}  
    & & & \multicolumn{2}{p{0.3\linewidth}}{
        \cellcolor[rgb]{0.9,0.9,0.9}{
            \makecell[{{p{\linewidth}}}]{
                \texttt{\tiny{[GM$|$GM]}}
                \texttt{yes} \\
            }
        }
    }
    & & \\ \\

    \theutterance \stepcounter{utterance}  
    & & & \multicolumn{2}{p{0.3\linewidth}}{
        \cellcolor[rgb]{0.9,0.9,0.9}{
            \makecell[{{p{\linewidth}}}]{
                \texttt{\tiny{[GM$|$GM]}}
                \texttt{game\_result = LOSE} \\
            }
        }
    }
    & & \\ \\

\end{supertabular}
}

\end{document}
